% ! TEX root = ../m3-19-20.tex

\chapter{Vectori aleatori bidimensionali}

Putem gîndi vectorii aleatori bidimensionali ca pe niște \emph{matrice %
  aleatoare}. Însă noțiunea este chiar ceva mai generală. Astfel, dacă
$ X $ și $ Y $ sînt două variabile aleatoare, atunci pur și simplu
formăm perechea $ (X, Y) $, care se numește \emph{vector aleator bidimensional}.
În majoritatea cazurilor, vom lucra cu $ X $ și $ Y $ definite pe același spațiu
de evenimente, deci dacă $ \dr{card}(X) = \dr{card}(Y) = n $, atunci perechea
lor devine o matrice aleatoare $ 2 \times n $.

\emph{Repartiția} unei astfel de matrice aleatoare se calculează intuitiv
folosind repartițiile individuale. Avem, deci:
\[
  P(X = x_i, Y = y_j) = p_{ij},
\]
unde se pot considera $ 1 \leq i \leq n $ și $ 1 \leq j \leq m $.
Mai notăm probabilitățile individuale cu $ p_i $, respectiv $ q_j $,
adică $ p_i = P(X = x_i) $ și $ q_j = P(Y = y_j) $. Aceste mărimi se mai
numesc \emph{probabilități marginale}.

Evident, dacă fixăm oricare dintre cele două variabile, obținem, de exemplu,
pentru $ X $ fixat:
\[
  \ds\cup_j (X = x_i, Y = y_j) = \sum_j P(X = x_i, Y = y_j) = \sum_j p_{ij} = p_i,
\]
deoarece ambele variabile aleatoare $ X $ și $ Y $ sînt, individual, normate.
Adică $ \sum_i p_i = \sum_i P(X = x_i) = \sum_j q_j = \sum_j P(Y = y_j) = 1 $.

Analog, desigur, $ \sum_i p_{ij} = q_j $.

Luate împreună, obținem automat condiția de normare pentru perechea de
variabile aleatoare, adică $ \sum_{i,j} p_{ij} = 1 $, în ipoteza că ambele
componente ale perechii sînt normate.

Notăm cu $ F_{(X, Y)}(x, y) = P(X < x, Y < y) $ funcția de repartiție
a perechii de variabile aleatoare (vectorului aleator bidimensional) $ (X, Y) $.
Separat, $ F_X(x) $ și $ F_Y(y) $ se numesc \emph{funcții de repartiție marginale}
ale vectorului aleator $ (X, Y) $.

Pentru \emph{cazul continuu}, întîlnim noțiunea cunoscută din cazul vectorilor
aleatori 1-dimensionali. Mai precis, avem:
\[
  F_{(X, Y)} (x, y) = \int_{-\infty}^x \int_{-\infty y} f(u, v) dudv,
\]
unde $ f $ este \emph{densitatea de probabilitate} și este o funcție cu
valori pozitive, dublu-normată, adică:
\[
  \int_{-\infty}^\infty \int_{-\infty}^\infty f(u, v) du dv = 1.
\]

Totodată, putem recupera din această ecuație și densitatea de probabilitate
știind funcția de repartiție, prin:
\[
  f(x, y) = \dfrac{\partial^2 F_{(X, Y)}(x, y)}{\partial x \partial y},
\]
evident, pentru punctele $ (x, y) $ unde egalitatea are sens.

%%%%%%%%%%%%%%%%%%%%%%%%%%%%%%%%%%%%%%%%%%%%%%%%%%%%%%%%%%%%%%%%%%%%%%
\section{Repartiții condiționate}

Pornim cu un vector aleator bidimensional $ (X, Y) $, discret.
Ne punem problema din cazul probabilităților condiționate, i.e. vom nota
cu $ P(x_i/y_j) $ probabilitatea să aibă loc evenimentul $ X = x_i $,
știind că a avut loc (i.e.\ condiționat de) evenimentul $ Y = y_j $.

Aceasta se poate calcula simplu, ca în cazul probabilităților condiționate:
\[
  P(x_i / y_j) = \dfrac{P(X = x_i, Y = y_j)}{P(Y = y_j)} = \dfrac{p_{ij}}{q_j}.
\]
Invers, obținem desigur $ P(y_j / x_i) = \dfrac{p_{ij}}{p_i} $.

Dacă fixăm un eveniment $ Y = y $, putem defini o \emph{funcție de repartiție %
  condiționată}, $ F(x/y) = P(X < x / Y = y) $, asociată vectorului $ X $
(deoarece $ Y $ este fixat). Aceasta se calculează generalizînd formulele de
mai sus pe cazul continuu, anume:
\[
  F(x/y) = \dfrac{1}{f_Y(y)} \cdot \int_{-\infty}^x f(u, v) du.
\]

Prin derivare în raport cu $ x $, obținem densitatea de probabilitate a
lui $ X $, în ipoteza că $ Y = y $:
\[
  f_X(x/y) = \dfrac{f(x, y)}{f_Y(y)}.
\]

%%%%%%%%%%%%%%%%%%%%%%%%%%%%%%%%%%%%%%%%%%%%%%%%%%%%%%%%%%%%%%%%%%%%%%
\section{Variabile aleatoare independente}

Știm deja că două variabile aleatoare $ X $ și $ Y $ sînt independente
dacă evenimentele $ X = x_i $ și $ Y = y_j $ sînt independente,
adică dacă are loc:
\[
  P(X = x_i, Y = y_j) = P(X = x_i) \cdot P(Y = y_j).
\]
Egalitatea poate fi formulată și cun funcțiile de repartiție:
\[
  F_{(X, Y)}(x, y) = F_X(x) \cdot F_Y(y).
\]

Mai mult decît atît, rezultatul de mai jos ne arată că egalitatea are
loc și pentru densitățile de probabilitate:

\begin{theorem}\label{thm:indep}
  Fie $ X, Y $ variabile aleatoare cu densitățile de probabilitate
  $ f_X $, respectiv $ f_Y $. Atunci $ X $ și $ Y $ sînt independente
  dacă și numai dacă are loc:
  \[
    f(x, y) = f_X(x) \cdot f_Y(y),
  \]
  unde $ f(x, y) $ este densitatea de probabilitate a vectorului $ (X, Y) $,
  format de cele două variabile aleatoare.
\end{theorem}

\textbf{Atenție:} Dacă lucrăm cu alte variabile aleatoare $ U, V $, obținute
din $ X $ și $ Y $ prin transformări funcționale de forma generală:
\[
  U = \varphi(X, Y), \quad V = \psi(X, Y),
\]
atunci calculele care folosesc densitățile de probabilitate și funcțiile
de repartiție se fac folosind schimbări de variabile de forma:
\[
  u = \varphi(x, y), v = \psi(x, y) \Leftrightarrow %
  x = \varphi_1(u, v), y = \psi_1(u,v),
\]
după care noua densitate de probabilitate se schimbă cu ajutorul jacobianului
transformărilor de mai sus:
\[
  f_{(U, V)}(u,v) = \left| \dfrac{D(x, Y)}{D(u,v)} \right| \cdot %
  f_{(X, Y)}(\varphi_1(u,v), \psi_1(u, v)),
\]
unde:
\[
  \left| \dfrac{D(x, y)}{D(u,v)} \right| = \det %
  \begin{pmatrix}
    x_u & x_v \\
    y_u & y_v
  \end{pmatrix},
\]
determinantul matricei jacobiene a transformării ($ x_u = \dfrac{\partial x}{\partial u} $
etc.).

\vspace{0.3cm}

Dintre toate caracteristicile numerice asociate variabilelor aleatoare
(medie, dispersie, coeficient de corelație, covarianță), pentru cazul
bidimensional avem în plus \emph{media condiționată}. Astfel, fie $ X, Y $
variabile aleatoare \emph{discrete}. Atunci putem calcula mediile condiționate
prin:
\begin{itemize}
\item $ E(X / Y = y_j) = \sum_i x_i \cdot P(x_i / y_j) $;
\item $ E(Y / X = x_i) = \sum_j y_j P(y_j / x_i) $.
\end{itemize}

Pentru \emph{cazul continuu}, sumele devin integrale, iar probabilitățile
devin densități:
\begin{itemize}
\item $ E(X / Y = y) = \int_{-\infty}^\infty x \cdot f_X(x/y) dx $;
\item $ E(Y / X = x) = \int_{-\infty}^\infty y \cdot f_Y(y/x) dy $.
\end{itemize}

%%%%%%%%%%%%%%%%%%%%%%%%%%%%%%%%%%%%%%%%%%%%%%%%%%%%%%%%%%%%%%%%%%%%%%
\section{Exerciții}

\indent\indent 1. Fie variabilele aleatoare:
\[
  X = \begin{pmatrix}
    -1 & 2 \\
    0.5 & 0.5
  \end{pmatrix}, \quad
  Y = \begin{pmatrix}
    -2 & 0 & 1 \\
    0.3 & 0.3 & 0.4
  \end{pmatrix}.
\]

\begin{enumerate}[(a)]
\item Calculați $ E(X / Y = 1) $;
\item Calculați $ D^2(X) $ și $ D^2(Y) $;
\item Calculați $ XY $ și apoi $ E(XY) $;
\item Determinați coeficientul de corelația $ \rho(X, Y) $.
\end{enumerate}

\vspace{0.3cm}

2. Fie vectorul aleator $ (X, Y) $, cu densitatea de repartiție:
\[
  f(x, y) = \begin{cases}
    e^{-y}, & 0 < x < y \\
    0, & \text{ în rest}
  \end{cases}.
\]
Determinați:
\begin{enumerate}[(a)]
\item Densitățile marginale $ f_X(x) $ și $ f_Y(y) $;
\item Densitatea variabilei aleatoare $ Y $ condiționată de $ X $;
\item Probabilitatea ca $ (x, y) \in (X, Y) $ să aparțină pătratului
  unitate $ [0, 1] \times [0, 1] $.
\end{enumerate}

\vspace{0.3cm}

3. Fie vectorul aleator $ (X, Y) $, cu densitatea de repartiție:
\[
  f(x, y) = \begin{cases}
    k, & 0 < y \leq x < 1 \\
    0, & \text{ în rest}
  \end{cases},
\]
unde $ k \in \RR $ este o constantă fixată.
\begin{enumerate}[(a)]
\item Determinați $ k $, astfel încît funcția de mai sus să fie corect
  definită, ca densitate de repartiție;
\item Determinați densitățile de repartiție marginale $ f_X(x), f_Y(y) $;
\item Calculați $ P(0 < X < 0.5, 0 < Y < 0.5) $;
\item Determinați densitățile condiționate $ f(y/x) $ și $ f(x/y) $.
\end{enumerate}

\vspace{0.3cm}

4. Fie funcția:
\[
  f(x, y) = \begin{cases}
    x^2 + \dfrac{xy}{c}, & x \in (0, 1), y \in (0,2) \\
    0, & \text{ în rest}
  \end{cases}.
\]
\begin{enumerate}[(a)]
\item Determinați $ c \in \RR $ astfel încît $ f(x, y) $ să fie o densitate
  de repartiție;
\item Determinați funcția de repartiție $ F_{(X,Y)}(x, y) $ a vectorului aleator
  bidimensional $ (X, Y) $ care are densitatea de repartiție calculată mai sus;
\item Determinați densitățile marginale $ f_X(x), f_Y(y) $ și funcțiile de
  repartiție $ F_X(x), F_Y(y) $;
\item Calculați $ P(X > 0.5) $ și $ P(Y < 0.5 / X < 0.5) $;
\item Calculați densitatea de repartiție condiționată $ F(Y | X = x) $.
\end{enumerate}

\vspace{0.3cm}

5. Vectorul aleator $ (X, Y) $ se dă prin tabloul de repartiție comună:
\[
  \begin{tabular}{l|l|l|l|l}
    X \mid Y & y_1 = -2 & y_2 = -1 & y_3 = 4 & y_4 = 5 \\
    \hline
    x_1 = 1 & p_{11} = 0.1 & p_{12} = 0.2 & p_{13} = 0 & p_{14} = 0.3 \\
    \hline
    x_2 = 2 & p_{21} = 0.2 & p_{22} = 0.1 & p_{23} = 0.1 & p_{24} = 0
  \end{tabular}.
\]
\begin{enumerate}[(a)]
\item Determinați repartițiile marginale $ p_X $ și $ p_Y $ ale celor două
  variabile aleatoare $ X $ și $ Y $;
\item Aflați dacă variabilele aleatoare sînt independente;
\item Calculați $ E(X) $ și $ E(Y) $;
\item Calculați $ \dr{cov}(X, Y) $ si determinați dacă variabilele aleatoare sînt
  corelate sau nu;
\item Calculați $ \rho(X, Y) $.
\end{enumerate}

\vspace{0.3cm}

6. Vectorul aleator $ (X, Y) $ se dă prin tabloul de repartiție comună:
\[
  \begin{tabular}{l|l|l}
    X \mid Y & y_1 = 0 & y_2 = 1 \\
    \hline
    x_1 = -1 & p_{11} = \lambda & p_{12} = 3/20 \\
    \hline
    x_2 = 0 & p_{21} = 3/20 & p_{22} = 1/4 \\
    \hline
    x_3 = 1 & p_{31} = 1/4 & p_{32} = 3/20
  \end{tabular}.
\]
\begin{enumerate}[(a)]
\item Determinați parametrul $ \lambda $;
\item Determinați $ D^2(X) $ și $ D^2(Y) $;
\item Verificați dacă cele două variabile aleatoare sînt independente.
\end{enumerate}

\vspace{0.3cm}

7. Vectorul aleator $ (X, Y) $ se dă prin tabloul de repartiție comună:
\[
  \begin{tabular}{l|l|l|l}
    X \mid Y & y_1 = 0 & y_2 = 1 & y_3 = 2 \\
    \hline
    x_1 = 2 & p_{11} = 0.2 & p_{12} = 0.3 & p_{13} = 0.1 \\
    \hline
    x_2 = 3 & p_{21} = 0.1 & p_{22} = 0.1 & p_{23} = \lambda
  \end{tabular}.
\]
\begin{enumerate}[(a)]
\item Determinați parametrul $ \lambda $;
\item Determinați repartițiile marginale ale vectorului aleator $ (X, Y) $;
\item Calculați $ E(Y / X = 2) $ și $ E(X / Y = 2) $.
\end{enumerate}

\vspace{0.3cm}

\textbf{Observație:} Exercițiile din seminar valorează fiecare cîte 2 puncte.



%%% Local Variables:
%%% mode: latex
%%% TeX-master: "../m3-19-20"
%%% End:
